\subsection{What the evidence shows, by dimension}
The findings concern different properties of the artifact, which should not be read as evidence for one another (Table~\ref{tab:dims}). Feasibility and solution quality are established on synthetic instances; computational efficiency is mixed; service reliability depends on the planning model more than on the search; managerial usefulness has not been tested.

\begin{table}[!htbp]
\centering\footnotesize
\caption{Findings by evaluation dimension.}\label{tab:dims}
\begin{tabular}{p{0.17\linewidth}p{0.47\linewidth}p{0.27\linewidth}}
\toprule
Dimension & Finding & Evidence and scope \\
\midrule
Feasibility & generic penalty methods rarely feasible under binding resources; all adapted methods always feasible & E1, E2, X2, X3, X5; synthetic \\
Solution quality & FMCG-ABC cheapest on the main instance and most instance groups; within 0.1\% of a certified optimum without discounts; about 8\% median and up to 47\% above the exact optimum with discounts for all adapted methods & E1--E3, X1--X5; certified or exact references where stated \\
Computational efficiency & 10--20 times slower per evaluation than vectorised GA+/PSO+; at equal time lowest median, significantly better than GA+ but not significantly better than PSO+ & E1, X1; one instance, one implementation \\
Service reliability & model-feasible policies miss individual targets for about half of the SKUs under daily review; the review-period variant reduces this to 9 of 50 at 3.9\% more cost & E7, S3, X6; synthetic holdout \\
Managerial usefulness & not evaluated & proposed study (Section~\ref{sec:limits}) \\
\bottomrule
\end{tabular}
\end{table}

The gap measures need the same care. Without discounts the bound is tight, so a gap to the bound is the true optimality gap. With discounts the bound can lie well below the optimum (the best known solution of the main instance is 18.7\% above it), so a large gap to the bound may reflect the bound rather than the policy; for the small instances, where the optimum is known, the methods themselves were about 8\% above it in the median. Planners should therefore read the reported gap as a guarantee (the policy is at most that far from optimal) and not as an estimate of the actual distance.

\subsection{Answers to the research questions}
\emph{RQ1.} The joint model represents the cost and service trade-offs of order quantity and safety stock with an expected-shortage cost and couples SKUs through the budget and the warehouses. Proposition~\ref{prop:convex} shows that its only source of non-convexity is the choice of price regimes, and the experiments confirm that this is where solution methods fail: without discounts, the adapted methods came within 0.1--0.2\% of a certified optimum, whereas on small discounted instances most runs selected a suboptimal combination of price regimes and stopped a median of about 8\%, and at worst 47\%, above the optimum.

\emph{RQ2.} The adaptations that matter most are inventory-specific constraint handling (service as a bound, repair, feasibility-first ranking), which turned a search that almost never returned a feasible policy into one that always did, and the log-scale multi-dimensional move, which reduced the median gap to the best known solution from 82\% to 15\%. EOQ-informed initialisation helped but did not, on its own, make the penalised search reliable. The constraint handling transfers to GA and PSO, so it is a property of the problem encoding rather than of the bee-colony metaphor.

\emph{RQ3.} FMCG-ABC returned the cheapest feasible policies on the main instance, with large and statistically significant differences at an equal evaluation budget. When the adapted GA and PSO received equal computing time, and therefore 10 to 19 times more evaluations, its median remained lowest; the difference was significant against GA+ but not against PSO+ with ten runs per method. It had the lowest median gap on the robustness instances, for portfolios of 50 to 100 SKUs over three instances per size, and under the average-investment budget. It was not the best method for 25-SKU portfolios, where PSO+ led, and on the small discounted instances the three adapted methods were practically tied. We therefore claim an advantage on these synthetic instance families, not general superiority. The more robust finding is negative and practical: generic penalty-based metaheuristics, as often used in inventory applications, returned inadmissible policies under binding resources.

\emph{RQ4.} The policies remained feasible under tighter and looser resources and for larger portfolios. The policy cost was most sensitive to the warehouse capacity, the budget and the service targets and little affected by demand variability. Out of sample, however, model-feasible policies missed individual service targets for about half of the SKUs, while fill rates stayed above 99\%. The diagnostic replay attributes most of the shortfall to the continuous-review assumption (undershoot of the reorder point under daily review) rather than to the normal approximation of demand, and building the review period into the model reduced the misses from 27 to 9 SKUs at 3.9\% higher controllable cost. This was a limitation of the planning model, not of the search, and it is the kind of discrepancy that an optimisation-only system would hide.

\emph{RQ5.} The artifact turns a policy into statements a planner can check: which cost component drives each SKU, which resources bind, how far the policy is from a lower bound, where model-implied and realised service diverge, and which rule fired on which evidence. The evaluation also showed that resource rules must be stated at portfolio level to remain informative.

\subsection{Implications for decision support design}
Three design principles generalise beyond inventory. First, \emph{feasibility is a first-class output}: a DSS that ranks candidate decisions by a penalised objective can recommend decisions that violate hard business limits by a small normalised amount but a large absolute one. Encoding the constraints the decision maker cannot negotiate as bounds or repairs, and reporting infeasibility rather than the least-bad candidate, prevented this. Second, \emph{a heuristic recommendation should carry a certificate}: a lower bound computed on the same instance turns ``the algorithm found this'' into ``this policy is at most $x$\% more expensive than any feasible policy'', which is an explanation of solution quality that managers and auditors can use, and in the convex case it is exact. Third, \emph{explanations must include what the model does not know}: reporting realised next to model-implied service exposed a systematic optimism of the planning model that an explanation of the optimisation alone would not reveal.

\subsection{Comparison with prior work}
Compared with the closest studies (Table~\ref{tab:position}), the improved ABC of \citet{Guo2025} addresses a single-item location--inventory problem with constant lead times and reports managerial insight through sensitivity analysis, and the decision support system of \citet{Long2023} optimises dynamic prices for a deteriorating product; neither treats several items competing for budget and storage, quantity discounts, or the quality of the returned solution. The multi-item studies closest in problem setting \citep{Pasandideh2013,Tan2024} treat discounts or a budget together with several items and compare metaheuristics, but report neither a bound on solution quality nor out-of-sample service, and neither embeds the result in an explanation layer. Classical multi-item results \citep{HadleyWhitin1963,RosenblattRothblum1990} provide the Lagrangian structure that our bound exploits but not the discounts, the service constraints or the decision support layer. Relative to the ABC literature \citep{Karaboga2008,Akay2012,Kaya2022}, our evidence is that, on this problem class, the problem-specific constraint handling matters more than the choice of swarm metaphor: the same handling made GA and PSO competitive. Relative to the explainable-AI literature for operational research \citep{DeBock2024,Coussement2024}, the artifact explains a decision rather than a prediction and attaches a bound on its quality, but, unlike the evaluation those authors call for, its effect on decision quality has not yet been measured.

\subsection{Managerial implications}
For FMCG planners, the results show that planning each SKU independently understates the controllable cost of a feasible policy by about half on the studied portfolio, that the budget and warehouse limits, not demand variability, drive the policy, and that service targets set in a continuous-review model should be checked against the review practice actually used. On the main instance, the multipliers of the bound put the value of an additional currency unit of procurement budget at about 0.9 currency units of annual controllable cost and identify the second warehouse as the binding space, which tells a planner where additional resources would pay off before any policy is changed. The DSS supports these checks without taking the decision: recommendations are accepted, modified or rejected by the planner, and the decision and its evidence are recorded.

\subsection{Limitations and future research}\label{sec:limits}
\emph{Data.} All results are synthetic. The generator is documented and its parameters are plausible for FMCG distribution, but external validity requires company data. The 30-SKU distributor dataset of the first author's thesis is supported by the IID research schema (Table~\ref{tab:dict}), but before it can be used, the period of its demand standard deviation, the units of lead time and volume, the basis of its holding and shortage costs, the budget, the warehouse capacities and assignment, the service targets and the missing lead-time variability and discount data must be confirmed or justified.

\emph{Model.} The main model limits the value of one batch of every SKU; the average-investment form changes cost levels but not the comparison (X5), and which form applies must be decided for each firm. The model assumes independent normal lead-time demand and backorders; the daily-review variant (X6) removes most of the service shortfall, but 9 SKUs still missed their targets, which a full periodic-review formulation or simulation-based service constraints could address. The warehouse constraint uses average stored volume; peak-volume constraints, case packs and minimum order quantities can be added without changing the method.

\emph{Method and evaluation.} Our ABC implementation is not vectorised, so its absolute runtime is 10 to 20 times that of GA+ and PSO+ per evaluation; at equal time it kept the lowest median on the main instance, but its advantage over PSO+ was not significant with ten runs per method; a vectorised implementation and larger equal-time comparisons remain open. The instance families are generated from one documented generator, and all methods use parameters fixed in advance from the literature rather than tuned per method, which treats them equally but does not give each its best configuration. The bound is not tight with discounts, so on discounted instances the reported gaps are upper bounds; branching over price regimes guided by the dual would tighten it, and the regime choice is where all methods lose most.

\emph{Usefulness to decision makers.} The explanation layer has been evaluated for fidelity, not for its effect on decisions. A feasible next study is a within-subject experiment with inventory planners or supply chain students (for example 20--30 participants) who make replenishment and budget-request decisions on several synthetic portfolios with (i) the optimised policy alone and (ii) the full artifact, in randomised order. Pre-registered outcomes would be the cost and service of the final decisions under replay, the detection rate of deliberately seeded errors in the recommendations, decision time, the frequency and justification of overrides recorded by the IID workflow, and perceived usefulness and ease of use. Until such a study is done, the DSS contribution of this paper is a designed, implemented and technically evaluated artifact, not evidence of better decisions.
